\markdownRendererDocumentBegin
\markdownRendererWarning{The `hybrid` option has been soft-deprecated.}{The `hybrid` option has been soft-deprecated.}{Consider using one of the following better options for mixing TeX and markdown: `contentBlocks`, `rawAttribute`, `texComments`, `texMathDollars`, `texMathSingleBackslash`, and `texMathDoubleBackslash`. For more information, see the user manual at <https://witiko.github.io/markdown/>.}{Consider using one of the following better options for mixing TeX and markdown: `contentBlocks`, `rawAttribute`, `texComments`, `texMathDollars`, `texMathSingleBackslash`, and `texMathDoubleBackslash`. For more information, see the user manual at <https://witiko.github.io/markdown/>.}\markdownRendererSectionBegin
\markdownRendererSectionBegin
\markdownRendererHeadingTwo{Patrones y Diseño de Algoritmos}\markdownRendererInterblockSeparator
{}\markdownRendererSectionBegin
\markdownRendererHeadingThree{Hamming Distance}\markdownRendererParagraphSeparator
{}La distancia de Hamming entre dos cadenas $a$ y $b$ de igual longitud es el número de posiciones donde difieren. Por ejemplo, $hamming(01101, 11001) = 2$.\markdownRendererParagraphSeparator
{}Si representamos las cadenas como enteros (cuando $k \le 32$), podemos calcular la distancia usando operaciones a nivel de bit. El XOR ($\oplus$) identifica los bits diferentes y contar los 1s.\markdownRendererParagraphSeparator
{}Esto reduce la complejidad de O($n^2 k$) a O($n^2$) al comparar $n$ cadenas de longitud $k$.\markdownRendererInterblockSeparator
{}\markdownRendererInputFencedCode{./_markdown_main/8582959398df3da9526e553164206b82.verbatim}{cpp}{cpp}\markdownRendererInterblockSeparator
{}$\clearpage$\markdownRendererInterblockSeparator
{}
\markdownRendererSectionEnd \markdownRendererSectionBegin
\markdownRendererHeadingThree{Reachability in Graphs}\markdownRendererInterblockSeparator
{}Dado un DAG con $n$ nodos, queremos calcular $reach(x)$: número de nodos alcanzables desde $x$.\markdownRendererParagraphSeparator
{}Se puede usar programación dinámica clásica O($n^2$) o bit-paralela usando \markdownRendererCodeSpan{bitset} para cada nodo, donde \markdownRendererCodeSpan{reach[x] \markdownRendererPipe{}= reach[u]} para cada $u$ adyacente a $x$. Esto reduce operaciones por nodo usando OR de bits, pero puede usar mucha memoria ($O(n^2)$ bits).\markdownRendererInterblockSeparator
{}\markdownRendererInputFencedCode{./_markdown_main/101b1ad853a3fc625f34f6f262366a71.verbatim}{cpp}{cpp}\markdownRendererInterblockSeparator
{}
\markdownRendererSectionEnd \markdownRendererSectionBegin
\markdownRendererHeadingThree{Counting Subgrids}\markdownRendererInterblockSeparator
{}Dado un grid $n \times n$ con casillas negras (1) o blancas (0), queremos contar subgrids cuyos cuatro vértices son negros.\markdownRendererParagraphSeparator
{}En la implementación bit-paralela, cada fila se representa como un \markdownRendererCodeSpan{bitset} y usamos operaciones AND para contar columnas donde ambas filas tienen negro. La complejidad baja de O($n^3$) a O($n^2$) y es mucho más rápida, pero usa más memoria ($O(n^2)$ bits para los bitsets).\markdownRendererInterblockSeparator
{}\markdownRendererInputFencedCode{./_markdown_main/efd015706bd2c9c5f243ce93c03cf8fb.verbatim}{cpp}{cpp}\markdownRendererInterblockSeparator
{}$\clearpage$\markdownRendererInterblockSeparator
{}
\markdownRendererSectionEnd \markdownRendererSectionBegin
\markdownRendererHeadingThree{Two Pointers Method}\markdownRendererInterblockSeparator
{}Se usan dos punteros $l$ y $r$ que solo avanzan hacia la derecha para encontrar subarrays con ciertas propiedades (ej: suma = x).\markdownRendererParagraphSeparator
{}Ejemplo clásico: encontrar subarray con suma exacta x. Ambos punteros se mueven O($n$) pasos en total, dando complejidad O($n$).\markdownRendererParagraphSeparator
{}También se aplica al problema 2SUM: encontrar dos valores cuya suma sea x. Primero se ordena el array O($n\log n$) y luego se usan dos punteros desde extremos opuestos.\markdownRendererInterblockSeparator
{}\markdownRendererInputFencedCode{./_markdown_main/509de053210700908a0ccb310219c5aa.verbatim}{cpp}{cpp}\markdownRendererInterblockSeparator
{}
\markdownRendererSectionEnd \markdownRendererSectionBegin
\markdownRendererHeadingThree{Nearest Smaller Element (Amortized Stack)}\markdownRendererInterblockSeparator
{}Para cada elemento de un array, encontramos el primer elemento más pequeño que lo precede usando un stack.\markdownRendererParagraphSeparator
{}Cada elemento se agrega exactamente una vez y se elimina a lo sumo una vez → O($n$) tiempo. Es un ejemplo de análisis amortizado: operaciones individuales varían, pero el total es lineal.\markdownRendererInterblockSeparator
{}\markdownRendererInputFencedCode{./_markdown_main/1a86719234dc2cbc9f1ec51d05d60072.verbatim}{cpp}{cpp}\markdownRendererInterblockSeparator
{}
\markdownRendererSectionEnd \markdownRendererSectionBegin
\markdownRendererHeadingThree{Sliding Window Minimum}\markdownRendererInterblockSeparator
{}Se mantiene una ventana de tamaño fijo moviéndose sobre un array, y queremos el mínimo de cada ventana.\markdownRendererParagraphSeparator
{}Se usa un deque: los elementos se agregan al final y se eliminan del final si son mayores al nuevo elemento. El primero siempre es el mínimo. Cada elemento se agrega y elimina a lo sumo una vez → O($n$).\markdownRendererParagraphSeparator
{}\markdownRendererInputFencedCode{./_markdown_main/4edc1587ad92a9bbbf6d818bef3b7e4e.verbatim}{cpp}{cpp}\markdownRendererInterblockSeparator
{}
\markdownRendererSectionEnd \markdownRendererSectionBegin
\markdownRendererHeadingThree{Minimizing Sums}\markdownRendererInterblockSeparator
{}Dado un conjunto de $n$ números $a_1, a_2, \dots, a_n$, consideremos el problema de encontrar un valor $x$ que minimice la suma\markdownRendererParagraphSeparator
{}$$|a_1 - x| + |a_2 - x| + \dots + |a_n - x|$$\markdownRendererParagraphSeparator
{}Por ejemplo, si los números son [1, 2, 9, 2, 6], la solución óptima es $x = 2$, lo que produce la suma\markdownRendererParagraphSeparator
{}$$|1-2| + |2-2| + |9-2| + |2-2| + |6-2| = 12$$\markdownRendererParagraphSeparator
{}Cada función $|a_k - x|$ es convexa, así que la suma también lo es. Podríamos usar búsqueda ternaria para encontrar el $x$ óptimo, pero hay una solución más simple: el valor óptimo de $x$ siempre es la mediana de los números, es decir, el elemento central después de ordenar. Por ejemplo, [1, 2, 9, 2, 6] se ordena como [1, 2, 2, 6, 9], y la mediana es 2.\markdownRendererParagraphSeparator
{}La mediana siempre es óptima, porque si $x$ es menor que la mediana, la suma disminuye al aumentar $x$, y si $x$ es mayor, la suma disminuye al disminuir $x$. Si $n$ es par y hay dos medianas, ambas medianas y todos los valores entre ellas son soluciones óptimas.\markdownRendererParagraphSeparator
{}Ahora consideremos minimizar la función\markdownRendererParagraphSeparator
{}$$(a_1 - x)^2 + (a_2 - x)^2 + \dots + (a_n - x)^2$$\markdownRendererParagraphSeparator
{}Por ejemplo, con [1, 2, 9, 2, 6], la mejor elección es $x = 4$, produciendo la suma\markdownRendererParagraphSeparator
{}$$(1-4)^2 + (2-4)^2 + (9-4)^2 + (2-4)^2 + (6-4)^2 = 46$$\markdownRendererParagraphSeparator
{}Esta función también es convexa. Aunque podríamos usar búsqueda ternaria, hay una solución directa: el valor óptimo de $x$ es el promedio de los números. En este ejemplo, el promedio es\markdownRendererParagraphSeparator
{}$$(1 + 2 + 9 + 2 + 6)/5 = 4$$\markdownRendererParagraphSeparator
{}Esto se puede demostrar reescribiendo la suma como\markdownRendererParagraphSeparator
{}$$n x^2 - 2x(a_1 + a_2 + \dots + a_n) + (a_1^2 + a_2^2 + \dots + a_n^2)$$\markdownRendererParagraphSeparator
{}La última parte no depende de $x$, y la función restante forma una parábola con mínimo en $x = s/n$, donde $s = a_1 + a_2 + \dots + a_n$. Por lo tanto, el mínimo se alcanza en el promedio de los números $a_1, a_2, \dots, a_n$.\markdownRendererInterblockSeparator
{}\markdownRendererInputFencedCode{./_markdown_main/61b6e729390ee22d3ad135987fbffa2e.verbatim}{cpp}{cpp}
\markdownRendererSectionEnd 
\markdownRendererSectionEnd 
\markdownRendererSectionEnd \markdownRendererDocumentEnd